% hw4-set-relation.tex (slides)

\begin{frame}[allowframebreaks]{作业 4：集合与关系}

\begin{problem}[关系的运算]
  请证明
  \[(X \cap Y) \circ Z \subseteq (X \circ Z) \cap (Y \circ Z),\]
  并举例说明不能换成等号。
\end{problem}

\pause
\begin{slidesolution}
  \begin{itemize}
    \item 证明包含关系：取任意 $(a,c)\in (X\cap Y)\circ Z$，存在 $b$ 使 $(a,b)\in X\cap Y$ 与 $(b,c)\in Z$，因此 $(a,c)\in X\circ Z$ 且 $(a,c)\in Y\circ Z$，从而属于交集。
    \pause
    \item 反例示意：取 $X=\{(2,4)\},\; Y=\{(3,4)\},\; Z=\{(1,2),(1,3)\}$，可检验左侧为空而右侧非空，故不可换成等号。
  \end{itemize}
\end{slidesolution}
\handoutonly{
\begin{proof}
  取任意 $(a,c)$, 由定义可推导出包含关系, 并给出反例
  \(X=\{(2,4)\},\; Y=\{(3,4)\},\; Z=\{(1,2),(1,3)\}\).
\end{proof}
}

\pause
\begin{problem}[关系的性质]
  请证明 $R$ 对称且传递当且仅当 $R = R^{-1} \circ R$。
\end{problem}

\begin{slidesolution}
  \begin{itemize}
    \item 要证等价：若 $R$ 对称且传递，则任取 $(x,y)\in R$ 有 $(y,x)\in R$，再由传递得到 $R\subseteq R^{-1}\circ R$；反向也类似，分两步验证包含关系即可。
    \pause
    \item 关键点：把证明拆成两部分（$\subseteq$ 与 $\supseteq$），分别使用对称性与传递性来推导元素归属。
  \end{itemize}
\end{slidesolution}
\handoutonly{
\begin{proof}
  见讲义中的分步证明, 首先从对称+传递推出等式, 再反向证明等价性。
\end{proof}
}

\end{frame}
